\textbf{15.}
Определим $\Sigma_k\mathsf{3SAT}, \Pi_k\mathsf{3SAT}, \Sigma_k\mathsf{3DNFSAT}$ и $\Pi_k\mathsf{3DNFSAT}$ аналогично $\Sigma_k \mathsf{SAT}$ и $\Pi_k \mathsf{SAT}$ для формул в формате 3-КНФ и 3-ДНФ соответственно. Для каждого $k$ классифицируйте эти языки по уровням полиномиальной иерархии и докажите их полноту.

\textbf{Ответ.} 
$\Sigma_k\mathsf{3SAT}$ будет $\Sigma_k^p$-полным при нечётных $k$, а  $\Pi_k\mathsf{3SAT}$ — $\Pi_k^p$-
полным при чётных. Иначе они будут полными в $\Sigma_{k-1}^p$и $\Pi_{k-1}^p$, соотвественно. $\Pi_k\mathsf{3DNFSAT}$ будет $\Sigma_k^p$-полным при нечётных $k$, а  $\Sigma_k\mathsf{3DNFSAT}$ — $\Pi_k^p$-
полным при чётных. Иначе они будут полными в $\Sigma_{k-1}^p$и $\Pi_{k-1}^p$, соотвественно.

\textbf{Решение.} 
Для начала заметим, что языки $\overline{\Sigma_k\mathsf{3SAT}}$ и $\Pi_k\mathsf{3DNFSAT}$ сводятся друг к другу, поскольку 
$$
\neg \exists x_{1,1}, x_{1, 2}, \ldots x_{1, n_1} \forall x_{2, 1}, \ldots, x_{2, n_2} \ldots Q x_{k, 1}, \ldots,x_{k, n_k}\colon A \Longleftrightarrow$$ $$ \Longleftrightarrow \forall x_{1,1}, x_{1, 2}, \ldots x_{1, n_1} \exists x_{2, 1}, \ldots, x_{2, n_2} \ldots P x_{k, 1}, \ldots,x_{k, n_k}\colon \neg A
$$
где $Q$ --- какой-то квантор (в зависимости от четности $k$), $P$ --- отличный от $Q$ квантор, $A$ --- формула в 3-КНФ, зависящая от переменных $x_{1, 1}, \ldots x_{1, n_1}, \ldots x_{k, 1} \ldots x_{k, n_k}$, а формула $\neg A$, соответственно, формула в 3-ДНФ.
Аналогично, друг к другу сводятся $\overline{\Pi_k\mathsf{3SAT}}$ и $\Sigma_k\mathsf{3DNFSAT}$. Таким образом, достаточно решить задачу для $\Sigma_k\mathsf{3SAT}$ и $\Pi_k\mathsf{3SAT}$.

Во-первых, в силу построения, $\Sigma_{2k+1}\mathsf{3SAT}$ лежит в 
$\Sigma_{2k+1}^p$. Чтобы доказать полноту, сведем к нему $\Sigma_{2k+1}\mathsf{SAT}$. Для этого вспомним, как \sf SAT \rm сводится к \sf 3-SAT \rm. Это происходит путем добавления к исходным переменным формулы дополнительных переменных $y_1, \ldots, y_m$. Аналогично бескванторному случаю, мы заведем дополнительные переменные и свяжем их последним квантором. Поскольку всего кванторов нечетное число, это будет именно квантор существования, то есть, сведение работает.

То же самое работает и для $\Pi_{2k}\mathsf{3SAT}$, поскольку там тоже последний квантор --- это $\exists$.

Заметим, что если $A$ --- формула в 3-КНФ, зависящая от переменных $x_1, \ldots x_n$, и значения переменных $x_1, \ldots, x_t$ зафиксированы, то условие $\forall x_{t+1} \ldots x_n A(x_1 \ldots x_n)$ проверяется за полиномиальное время. Приведем алгоритм такой проверки. Рассмотрим множество дизъюнктов, входящих в $A$. Уберем все дизъюнкты, в которых есть литерал от какой-то переменной из $x_1, \ldots, x_t$, равный $1$. Если дизъюнктов больше не осталось, условие выполнено. Если среди оставшихся нашелся дизъюнкт, зависящий только от переменных $x_1, \ldots x_t$ и равный 0, условие не выполнено. Иначе условие будет выполнено тогда и только тогда, когда найдется дизъюнкт, содержащий одновременно $p$ и $\neg p$, что проверяется за полиномиальной время. 

Теперь докажем, что $\Sigma_{2k}\mathsf{3SAT}$ полон в $\Sigma_{2k-1}^p$. В силу предыдущего замечания, он лежит в этом классе, поскольку последний квантор - это $\forall$. А еще к нему можно свести язык $\Sigma_{2k-1}\mathsf{3SAT}$, полный в $\Sigma_{2k-1}^p$, просто добавив квантор $\forall$ по фиктивной переменной. Значит, $\Sigma_{2k}\mathsf{3SAT}$ полон в $\Sigma_{2k-1}^p$. Для $\Pi_{2k+1}\mathsf{3SAT}$ доказательство полноты в $\Pi_{2k}^p$ аналогичное.


